\chapter{AI Safety and Regulation}
\label{ch:ch89_ai_safety}

\section*{Patient focus}
If an AI tool helps your doctor interpret a scan, you might wonder: Is it safe? Has it been tested? Who is responsible if it makes a mistake? These are important questions. AI tools used in medicine go through rigorous testing and regulatory review before they reach your hospital. ``FDA cleared'' means the U.S. Food and Drug Administration (FDA) has evaluated the tool for safety and effectiveness for a specific purpose. However, AI is not perfect---it can sometimes make errors, just like any tool. That is why your doctor always reviews the final result. AI is there to assist, not to replace, your doctor's judgment.

\section*{Imaging focus}
AI safety in imaging encompasses regulatory pathways (510(k), De Novo, Premarket Approval [PMA]), bias and fairness, transparency, liability, human oversight, and post-market surveillance. Understanding how AI tools are regulated and where they can fail is essential for responsible clinical deployment.

\section{The Regulatory Landscape}
In the United States, the FDA regulates artificial intelligence and machine learning (AI/ML)-based software as a medical device (SaMD). Three primary pathways exist:

\begin{table}[ht]
\centering
\caption{FDA regulatory pathways for AI/ML medical devices}
\label{tab:fda_ai_pathways}
\begin{tabular}{@{}p{0.15\textwidth}p{0.42\textwidth}p{0.35\textwidth}@{}}
\toprule
\textbf{Pathway} & \textbf{Requirements} & \textbf{Example} \\
\midrule
510(k) & Substantial equivalence to a predicate device & AI triage for head CT \\
De Novo & Novel device with low-to-moderate risk, no predicate & Autonomous diabetic retinopathy screening \\
PMA & High-risk device requiring clinical trial evidence & AI-guided cardiac imaging \\
\bottomrule
\end{tabular}
\end{table}

A 510(k) clearance means the manufacturer demonstrated that their AI tool performs at least as well as an already-cleared device for the same intended use. It does not mean the tool is ``best in class'' or universally applicable.

\section{What ``FDA Cleared'' Does and Does Not Mean}
\begin{itemize}
    \item \textbf{Does mean}---the device was tested for a specific intended use, on representative data, with defined performance benchmarks
    \item \textbf{Does not mean}---the device works equally well on all patient populations, all scanner types, or all clinical settings
    \item \textbf{Does not mean}---the device is approved for uses beyond its labeled indications
\end{itemize}

\begin{pearlbox}[Intended use matters]
An AI tool cleared for detecting intracranial hemorrhage on head CT is not automatically cleared for detecting stroke, tumor, or fracture---even if the same images are used. Off-label use of AI tools is a regulatory and patient safety concern.
\end{pearlbox}

\section{Bias and Fairness}
AI models learn patterns from training data. If that data is not representative, the model may perform differently across demographic groups:

\begin{itemize}
    \item \textbf{Demographic bias}---underperformance on underrepresented groups (e.g., skin color in dermatology AI, age groups in pediatric imaging)
    \item \textbf{Selection bias}---training data from a single tertiary center may not reflect community hospital presentations
    \item \textbf{Label bias}---if ground-truth labels are influenced by biased clinical practices, the model inherits those biases
    \item \textbf{Measurement bias}---differences in protocol, scanner, or acquisition across sites
\end{itemize}

Fairness auditing requires evaluating model performance across sex, age, ethnicity, body habitus, and disease severity subgroups before deployment.

\section{Transparency and Explainability}
Clinicians and patients need to understand how an AI tool reaches its conclusions:

\begin{table}[ht]
\centering
\caption{Levels of AI transparency}
\label{tab:ai_transparency}
\begin{tabular}{@{}p{0.33\textwidth}p{0.58\textwidth}@{}}
\toprule
\textbf{Level} & \textbf{Description} \\
\midrule
Full interpretability & Human-understandable model internals \\
Model documentation & Architecture, training data, intended use, known limitations \\
Output explanation & Highlighted regions, attention maps, confidence scores \\
Reasoning trace & Step-by-step logic connecting input to output \\
\bottomrule
\end{tabular}
\end{table}

Most current medical AI provides output explanation (e.g., heatmaps showing regions of interest) rather than full interpretability. Clinicians should interpret these with caution---attention maps show where the model looked, not why it concluded what it did.

\section{Liability and Accountability}
When an AI tool contributes to a clinical decision, liability questions arise:
\begin{enumerate}
    \item \textbf{Manufacturer}---responsible for device safety and performance within labeled indications
    \item \textbf{Institution}---responsible for appropriate deployment, validation, and monitoring
    \item \textbf{Clinician}---responsible for the final clinical decision, regardless of AI output
\end{enumerate}

The reporting clinician always bears ultimate responsibility. ``The AI missed it'' is not a defence if the clinician failed to exercise independent judgment.

\begin{pearlbox}[The human-in-the-loop principle]
No AI tool should make autonomous clinical decisions in imaging. Every AI output must be reviewed by a qualified clinician who retains responsibility for the final interpretation and patient management.
\end{pearlbox}

\section{Human Oversight Requirements}
Safe AI deployment in radiology requires:
\begin{itemize}
    \item Display of original images alongside AI outputs
    \item Ability to override or dismiss AI recommendations
    \item Clear documentation of AI involvement in the report
    \item Training for radiologists on AI capabilities and limitations
    \item Defined pathways for reporting AI errors and near-misses
\end{itemize}

\section{Post-Market Surveillance}
Regulatory clearance is not the end of evaluation. Post-market activities include:
\begin{itemize}
    \item \textbf{Performance monitoring}---tracking sensitivity, specificity, and error rates over time
    \item \textbf{Drift detection}---identifying when model performance degrades due to changing patient populations or scanner upgrades
    \item \textbf{Adverse event reporting}---mandatory reporting of device-related patient harm
    \item \textbf{Algorithm updates}---FDA requires submission for changes that alter intended use or performance, unless pre-authorized under a Predetermined Change Control Plan (PCCP)
\end{itemize}

\section{When AI Can Harm}
AI can cause patient harm through:
\begin{enumerate}
    \item \textbf{False negatives}---missing a critical finding (e.g., pneumothorax, fracture) leading to delayed treatment
    \item \textbf{False positives}---triggering unnecessary workups, biopsies, or patient anxiety
    \item \textbf{Automation bias}---clinician over-reliance on AI, leading to inattention to the images
    \item \textbf{Data leakage}---patient information exposed through model inversion or membership inference attacks
    \item \textbf{Workflow disruption}---alert fatigue from excessive low-specificity notifications
\end{enumerate}

\section{Global Regulatory Comparison}
\begin{table}[ht]
\centering
\caption{Regulatory frameworks for medical AI}
\label{tab:ai_global_regulation}
\begin{tabular}{@{}p{0.25\textwidth}p{0.66\textwidth}@{}}
\toprule
\textbf{Region} & \textbf{Regulatory body / framework} \\
\midrule
China & National Medical Products Administration (NMPA) registration \\
European Union & CE marking under Medical Device Regulation (MDR 2017/745) and the EU AI Act \\
India & Central Drugs Standard Control Organisation (CDSCO) approval \\
United Kingdom & UK Conformity Assessed (UKCA) marking via the Medicines and Healthcare products Regulatory Agency (MHRA) \\
United States & FDA (510(k), De Novo, PMA, and PCCP) \\
\bottomrule
\end{tabular}
\end{table}

\section{Summary}
AI safety is a lifecycle property---it begins with data governance, continues through validation and deployment, and requires ongoing surveillance. Regulatory clearance provides a baseline assurance but does not guarantee universal performance. Clinicians must understand intended use, recognize bias, maintain independent judgment, and report failures. Patient safety depends on the combination of well-regulated AI tools and well-trained, vigilant clinicians.

\clearpage
